\section{多项式函数}

形式多项式由系数序列定义，而代入后得到的是函数。两种对象之间存在自然联系，但能否由函数值恢复多项式，需要分别讨论无限域和有限域。

\begin{definition}
设$F$为域，$F^F$表示所有从$F$到$F$的函数组成的集合。对$f,g\in F^F$，逐点定义
\[
(f+g)(a)=f(a)+g(a),\qquad(fg)(a)=f(a)g(a).
\]
常值函数$a\mapsto c$记作$f_c$。再定义函数$s\in F^F$为$s(a)=a$。
\end{definition}

\begin{proposition}
$F^F$是交换幺环，$c\mapsto f_c$将$F$嵌入其中。子环$F[s]$中的元素称为$F$上的\textbf{多项式函数}；它们恰是形如
\[
a\longmapsto c_0+c_1a+\cdots+c_na^n
\]
的函数。
\end{proposition}
\begin{proof}
函数相等当且仅当每一点的取值相等。逐点使用$F$中的加法群性质、乘法结合律及分配律，就得到函数上的相应运算律。零元与单位元分别为$f_0,f_1$，负元为$a\mapsto-f(a)$，乘法也可交换。

常值函数的运算与其常数的运算一致；两个不同常数给出不同函数，故常值函数构成$F$的一个副本。最后，按生成子环的定义，$F[s]$的元素恰为有限和$\sum c_i s^i$；在$a$处取值即得所列公式。
\end{proof}

由一元多项式环的延拓定理，有满同态
\[
\sigma:F[x]\longrightarrow F[s],\qquad
\sum_i c_i x^i\longmapsto\sum_i c_i s^i.
\]
它把形式不定元$x$映为函数$s$。记其核为$I$，则
\[
I=\{f\in F[x]:f(a)=0\text{ 对所有 }a\in F\text{ 成立}\},
\qquad F[s]\cong F[x]/I.
\]
所以问题归结为：哪些多项式在整个域上恒等于零？

\begin{theorem}\label{p2:poly-functions-one}
若$F$为无限域，则$I=\{0\}$，从而$F[x]\cong F[s]$，每个多项式函数由唯一的形式多项式表示。

若$F$为$q$元有限域，则
\[
I=(x^q-x),\qquad F[s]\cong F[x]/(x^q-x).
\]
此时每个从$F$到$F$的函数都是多项式函数，并且有唯一的次数小于$q$的多项式代表元。
\end{theorem}
\begin{proof}
\textbf{无限域：}若非零$f\in I$，设$\deg f=n$，则$f$在无限多个点取零值。但根数定理说明非零$n$次多项式至多有$n$个根，矛盾。因此$I=0$，满同态$\sigma$为同构。

\textbf{有限域中的核：}每个$a\in F$满足$a^q=a$，所以$x^q-x\in I$。另一方面，
\[
\prod_{a\in F}(x-a)=x^q-x.
\]
为证明此式，左边是首一$q$次多项式，并且两边在全部$q$个点取值为零。两者之差次数小于$q$却有$q$个不同的根，只能是零多项式。

对$f\in I$，用$x^q-x$作带余除法，写成
\[
f=h(x)(x^q-x)+r(x),\qquad\deg r<q.
\]
代入任意$a\in F$，得到$r(a)=0$。因此$r$有$q$个根而次数小于$q$，只能为零。所以$f\in(x^q-x)$；结合反向包含关系，得$I=(x^q-x)$。

\textbf{所有函数的表示：}带余除法说明每个剩余类有唯一的次数小于$q$的代表元。这些多项式共有$q^q$个，因为$q$个系数各有$q$种选择。不同代表元给出不同函数，否则差属于$I$且次数小于$q$，只能为零。

从$q$元集合到自身的函数也恰有$q^q$个。因此多项式函数已包含全部函数，且每个函数的上述代表元唯一。若不限制次数，则可以加上任意$(x^q-x)$的倍式，表示便不再唯一。
\end{proof}

\begin{definition}
对正整数$n$，令$S=F^n$，$F^S$表示从$S$到$F$的全部函数，仍按逐点运算组成交换幺环。定义第$i$个\textbf{坐标函数}
\[
s_i(a_1,\ldots,a_n)=a_i\qquad(1\le i\le n).
\]
常值函数与$s_1,\ldots,s_n$生成的子环$F[s_1,\ldots,s_n]$称为$n$元多项式函数环。
\end{definition}

多元延拓定理给出满同态
\[
\sigma:F[x_1,\ldots,x_n]\to F[s_1,\ldots,s_n],
\qquad f\longmapsto f(s_1,\ldots,s_n).
\]
其核为在$F^n$上处处取零值的多项式组成的理想。

\begin{theorem}\label{p2:poly-functions-many}
若$F$为无限域，则$s_1,\ldots,s_n$在$F$上代数无关，且
\[
F[x_1,\ldots,x_n]\cong F[s_1,\ldots,s_n].
\]
若$F$为$q$元有限域，则
\[
\ker\sigma=(x_1^q-x_1,\ldots,x_n^q-x_n).
\]
每个$F^n\to F$的函数都是多项式函数，并且可唯一地由一个对每个变量次数均小于$q$的多项式表示。
\end{theorem}
\begin{proof}
\textbf{无限域：}对变量个数$n$归纳。$n=1$已经证明。设$f$在$F^n$上处处为零，将它按最后一个变量展开为
\[
f=\sum_{j=0}^m f_j(x_1,\ldots,x_{n-1})x_n^j.
\]
固定$a_1,\ldots,a_{n-1}\in F$，所得一元多项式
\[
\sum_{j=0}^m f_j(a_1,\ldots,a_{n-1})x_n^j
\]
对每个$x_n\in F$取零值。由一元结论，它是零多项式，故每个系数$f_j(a_1,\ldots,a_{n-1})$为零。固定的点任意，因此每个$f_j$在$F^{n-1}$上处处为零。由归纳假设，$f_j=0$，所以$f=0$。这证明了$\ker\sigma=0$，也证明坐标函数代数无关。

\textbf{有限域中的降次：}记
\[
J=(x_1^q-x_1,\ldots,x_n^q-x_n).
\]
由于每个$a\in F$满足$a^q=a$，各生成元在$F^n$上处处为零，所以$J\subseteq\ker\sigma$。模$J$计算时，若$e\ge q$，则
\[
x_i^e-x_i^{e-q+1}=x_i^{e-q}(x_i^q-x_i)\in J.
\]
因$q\ge2$，每次替换使指数降低$q-1$。对每个单项式的每个指数反复进行，最终都小于$q$。因此任一$f$均可写成
\[
f=\sum_{i=1}^n h_i(x_1,\ldots,x_n)(x_i^q-x_i)+r,
\]
其中$r$对每个变量的次数均小于$q$。

\textbf{降次后零函数只能由零多项式表示：}仍对$n$归纳。$n=1$由根数定理成立。对$n>1$，将这样的$r$写成
\[
r=\sum_{j=0}^{q-1}r_j(x_1,\ldots,x_{n-1})x_n^j.
\]
若$r$在$F^n$上处处为零，固定前$n-1$个坐标后，所得关于$x_n$的多项式次数小于$q$，却在全部$q$个点取零值，所以各系数$r_j(a_1,\ldots,a_{n-1})$均为零。每个$r_j$对前面的变量次数仍小于$q$，由归纳假设，$r_j=0$，从而$r=0$。

若$f\in\ker\sigma$，上述分解中的$r$也给出零函数，所以$r=0$，得到$f\in J$。这证明$\ker\sigma=J$。两个降次后的代表元若给出相同函数，其差也是同类的零函数，故差为零，证明了唯一性。

\textbf{计数得到全部函数：}指数都小于$q$的单项式共有$q^n$个，因而这样的多项式共有$q^{q^n}$个。$F^n$有$q^n$个点，从它到$F$的函数同样共有$q^{q^n}$个。不同降次多项式已给出不同函数，数量又相同，所以全部函数均已得到。
\end{proof}

因此，在无限域上可以由全部函数值唯一确定形式多项式；在有限域上，多项式函数对应的是模上述关系理想的剩余类。讨论系数、次数或因式分解时，应始终分清所说的是形式多项式，还是它所定义的函数。
